% Options for packages loaded elsewhere
\PassOptionsToPackage{unicode}{hyperref}
\PassOptionsToPackage{hyphens}{url}
\PassOptionsToPackage{dvipsnames,svgnames,x11names}{xcolor}
\documentclass[
  10pt,
  a4paper,
]{article}
\usepackage{xcolor}
\usepackage[margin=18mm]{geometry}
\usepackage{amsmath,amssymb}
\setcounter{secnumdepth}{-\maxdimen} % remove section numbering
\usepackage{iftex}
\ifPDFTeX
  \usepackage[T1]{fontenc}
  \usepackage[utf8]{inputenc}
  \usepackage{textcomp} % provide euro and other symbols
\else % if luatex or xetex
  \usepackage{unicode-math} % this also loads fontspec
  \defaultfontfeatures{Scale=MatchLowercase}
  \defaultfontfeatures[\rmfamily]{Ligatures=TeX,Scale=1}
\fi
\usepackage{lmodern}
\ifPDFTeX\else
  % xetex/luatex font selection
  \setmainfont[]{DejaVu Serif}
  \setsansfont[]{DejaVu Sans}
  \setmonofont[]{DejaVu Sans Mono}
\fi
% Use upquote if available, for straight quotes in verbatim environments
\IfFileExists{upquote.sty}{\usepackage{upquote}}{}
\IfFileExists{microtype.sty}{% use microtype if available
  \usepackage[]{microtype}
  \UseMicrotypeSet[protrusion]{basicmath} % disable protrusion for tt fonts
}{}
\makeatletter
\@ifundefined{KOMAClassName}{% if non-KOMA class
  \IfFileExists{parskip.sty}{%
    \usepackage{parskip}
  }{% else
    \setlength{\parindent}{0pt}
    \setlength{\parskip}{6pt plus 2pt minus 1pt}}
}{% if KOMA class
  \KOMAoptions{parskip=half}}
\makeatother
\usepackage{longtable,booktabs,array}
\newcounter{none} % for unnumbered tables
\usepackage{calc} % for calculating minipage widths
% Correct order of tables after \paragraph or \subparagraph
\usepackage{etoolbox}
\makeatletter
\patchcmd\longtable{\par}{\if@noskipsec\mbox{}\fi\par}{}{}
\makeatother
% Allow footnotes in longtable head/foot
\IfFileExists{footnotehyper.sty}{\usepackage{footnotehyper}}{\usepackage{footnote}}
\makesavenoteenv{longtable}
\usepackage{graphicx}
\makeatletter
\newsavebox\pandoc@box
\newcommand*\pandocbounded[1]{% scales image to fit in text height/width
  \sbox\pandoc@box{#1}%
  \Gscale@div\@tempa{\textheight}{\dimexpr\ht\pandoc@box+\dp\pandoc@box\relax}%
  \Gscale@div\@tempb{\linewidth}{\wd\pandoc@box}%
  \ifdim\@tempb\p@<\@tempa\p@\let\@tempa\@tempb\fi% select the smaller of both
  \ifdim\@tempa\p@<\p@\scalebox{\@tempa}{\usebox\pandoc@box}%
  \else\usebox{\pandoc@box}%
  \fi%
}
% Set default figure placement to htbp
\def\fps@figure{htbp}
\makeatother
\setlength{\emergencystretch}{3em} % prevent overfull lines
\providecommand{\tightlist}{%
  \setlength{\itemsep}{0pt}\setlength{\parskip}{0pt}}
\usepackage{xeCJK}
\setCJKmainfont[AutoFakeBold=1.4,AutoFakeSlant=0.2]{Droid Sans Fallback}
\setCJKsansfont[AutoFakeBold=1.4,AutoFakeSlant=0.2]{Droid Sans Fallback}
\setCJKmonofont{Droid Sans Fallback}
\emergencystretch=3em
\setlength{\parindent}{0pt}
\xeCJKDeclareCharClass{Default}{"2013,"00B7}
\usepackage{newunicodechar}
\newunicodechar{ℱ}{\ensuremath{\mathcal{F}}}
\renewcommand{\contentsname}{目录}

\usepackage{xurl}

\usepackage{float}
\floatplacement{figure}{H}
\renewcommand{\figurename}{图}
\usepackage{fancyhdr}
\pagestyle{fancy}
\fancyhf{}
\fancyhead[L]{\small 工作稿 — Watson 固定目标计算进行中}
\fancyfoot[C]{\thepage}
\renewcommand{\headrulewidth}{0.3pt}
\setlength{\headheight}{14pt}
\usepackage{bookmark}
\IfFileExists{xurl.sty}{\usepackage{xurl}}{} % add URL line breaks if available
\urlstyle{same}
\hypersetup{
  pdftitle={He--Kruczenski 2309.12402v3：ρ主线复现与原型机审计（工作稿）},
  colorlinks=true,
  linkcolor={Maroon},
  filecolor={Maroon},
  citecolor={Blue},
  urlcolor={Blue},
  pdfcreator={LaTeX via pandoc}}

\title{He--Kruczenski 2309.12402v3：ρ主线复现与原型机审计（工作稿）}
\usepackage{etoolbox}
\makeatletter
\providecommand{\subtitle}[1]{% add subtitle to \maketitle
  \apptocmd{\@title}{\par {\large #1 \par}}{}{}
}
\makeatother
\subtitle{Watson 固定目标计算进行中；非最终完成报告}
\author{}
\date{2026-09-11}

\begin{document}
\maketitle

{
\setcounter{tocdepth}{2}
\tableofcontents
}
2026-09-11。本文是本轮结果审计，旧报告及失败输入全部保留。当前仍有固定Watson目标计算在运行，其结果须在完成后更新本节，不能把本文工作稿当成全部复现完成的报告。完整数学论证另见\href{../proof_bundle_v2/PROOFS_ZH.pdf}{33页证明汇编}，运行索引见\href{../PROGRESS.json}{PROGRESS}。

\section{1.
核心判断：定性机制成立，正确且稳健的ρ尚未充分复现}\label{ux6838ux5fc3ux5224ux65adux5b9aux6027ux673aux5236ux6210ux7acbux6b63ux786eux4e14ux7a33ux5065ux7684ux3c1ux5c1aux672aux5145ux5206ux590dux73b0}

\href{../../../reports/expert_review_20260910/He_Kruczenski_2309v3_Expert_Review_ZH.pdf}{原专家报告第16页}的C3/Fig.7只有IR输入。该对照没有合适的ρ，并不是单凭这一页就能判定失败；论文真正要检验的是，在同一散射架构中加入UV电流、QCD谱矩和形状因子条件之后，预先选定的完整振幅是否出现合适、稳定的P波共振，并同时符合其他分波和分辨率比较。\href{https://arxiv.org/abs/2309.12402v3}{原论文§4、Fig.7--11}

本轮已完成IR→UV联合可行性、完整几何代表、相移/强度和声明网格上的一致核验。原生基线以及加强模型都出现了P1的90°上穿；完整振幅还给出可以严格括住的P1强度内部峰。但加强模型的三代表质量读数分散，P1存在明显非弹性，S0也有较大差异。因而不能把``tip接近770
MeV''或``能画出峰''称为论文核心定量结果已经成功。

这里没有证明原论文错误，也没有证明其全部允许域无法产生正确ρ。已证明的是本文具名有限完成的若干可行性、排除与峰值结论，以及已经实际完成的几何选择结果存在的定量差异。原作者没有公开的离散细节和代表点规则仍有识别缺口。

\section{2.
从底层约束到ρ：每一步到底推出什么}\label{ux4eceux5e95ux5c42ux7ea6ux675fux5230ux3c1ux6bcfux4e00ux6b65ux5230ux5e95ux63a8ux51faux4ec0ux4e48}

\subsection{2.1
同一个完整振幅}\label{ux540cux4e00ux4e2aux5b8cux6574ux632fux5e45}

共同解析族保留M=50时的3876个完整散射系数，包括所有cardinal模及Nyquist模。交叉对称、同位旋组合和部分波投影由同一个振幅生成，不能为不同图更换插值方式或单独重建P波。电流加入后有4076个原始变量，包括两个通道的ImF与谱R。

当前analytic-cardinal是一个明确的解析有限函数族；原PV/midpoint通路属于不同的混合配点近似。前者的秩、解析延拓、离节点求值和尾部结论不能移交后者，反之亦然。现已验证的是本文声明的完成方式，不是作者未公开程序逐字相同的恢复。\href{../../../evidence/ANALYTIC_CARDINAL_FINITE_PROOFS_20260911.md}{有限函数族证明}

散射归一化为

\[S_\ell^I=1+i\kappa f_\ell^I,\qquad S=\eta e^{2i\delta},\qquad |S|\le1.\]

低能条件以两个独立L2球施加，εχ=.002。它限制部分波之间的手征比例；单凭这些低能条件不会预置一个ρ质量。本文新的IR对照在所报能区末端只有δP1=8.07855°，S0零点严格括在无量纲s∈(.4,.5)。零点存在由两端严格异号及解析连续性推出，不声称唯一。

\subsection{2.2
UV通过电流Gram把散射与谱联结}\label{uvux901aux8fc7ux7535ux6d41gramux628aux6563ux5c04ux4e0eux8c31ux8054ux7ed3}

对于某一电流通道，设带运动学因子的形状因子为\(\mathcal F=kF\)、谱为R。Gram半正定包含\(R\ge|\mathcal F|^2\)。当R\textgreater0且F≠0时，令

\[r=|\mathcal F|^2/R,\qquad Q_F=F/F^*,\qquad 0\le r\le1.\]

在这些前提下，完整3×3条件精确等价于

\[|S-rQ_F|\le1-r.\]

F=0时不定义Q\_F，直接取\(rQ_F=\mathcal F^2/R=0\)。因此谱矩、F(0)=1及高能FF界会通过R、F限制S；这些约束不是给P波直接塞入一个Breit--Wigner峰。该圆盘还给出\(\eta\ge\max(0,2r-1)\)，但并不给出η≈1或r≈1。比如R=1、\(\mathcal F=i\sqrt{.6}\)、S=−.2在一个节点满足Gram半正定，却有δ=90°、η=.2。这个局部代数例子不被夸大成一条完整QCD反例；本轮完整联合振幅则实际展示了强非弹性也能通过声明的全部有限约束。\href{../../../evidence/GRAM_DISK_AND_INELASTICITY_20260911.md}{Gram证明}

\subsection{2.3
有90°读数，为什么仍可能没有``正确ρ''}\label{ux670990ux8bfbux6570ux4e3aux4ec0ux4e48ux4ecdux53efux80fdux6ca1ux6709ux6b63ux786eux3c1}

所画P1强度为

\[\mathcal I=|S-1|^2/4=(1+\eta^2-2\eta\cos2\delta)/4.\]

δ=90°时，\(\mathcal I=(1+\eta)^2/4\)。η≈1的共振和η≈.2时的90°点具有不同的散射强度和物理解释。若S靠近零，\(d\delta=\frac12\operatorname{Im}(dS/S)\)还说明相位读数更敏感。因此质量、强度、η以及其他分波必须一起比较，不能只对齐一个横坐标。

新三点均以严格区间值验证P1峰候选的中间值高于两侧；连续性保证括区内存在内部最大值。这个证明不依赖拟合论文曲线，但只证明存在峰，不证明唯一峰、精确质量或第二张黎曼面极点。FF的细小局部起伏也不逐个称为介子。\href{../FINE_SIGNAL_PEAK_PROOF_ZH.md}{新峰证明}

\subsection{2.4
已证UV机制与尚未证的更强命题}\label{ux5df2ux8bc1uvux673aux5236ux4e0eux5c1aux672aux8bc1ux7684ux66f4ux5f3aux547dux9898}

将新的完整IR振幅C固定，只让电流自由变化，在原2309的UV条件下得到解析计入误差的排除上界−.1484093366072121。它没有借用新增FF端点或高能散射条件，故确实说明这一个IR解不能直接延伸成原UV联合解。\href{../FINE_IR_UV_EXCLUSION.json}{排除证据}

另有底层推导证明有限UV条件强制至少一处P波有正吸收量，及足够弱的小振幅IR族受归一化和谱矩排除。但这些较弱定理没有推出所有可行解都具有同一质量的ρ，更不能把一个固定C的排除变成整个无ρ集合的排除。\href{../../../evidence/UV_FORCED_P_WAVE_ABSORPTION_20260911.md}{吸收下界}、\href{../../../evidence/GLOBAL_IR_SEED_AND_UV_OBSTRUCTION_20260911.md}{IR族与UV障碍}

\section{3.
本轮实际修补及没有改变的物理输入}\label{ux672cux8f6eux5b9eux9645ux4feeux8865ux53caux6ca1ux6709ux6539ux53d8ux7684ux7269ux7406ux8f93ux5165}

mπ=.140、fπ=.092 GeV，匹配能量1.2
GeV；Nc=3、Nf=2、αs=.4、mu=.004、md=.0073 GeV，凝聚量G=.023
GeV⁴、Q=−.0001
GeV⁴。轻夸克质量采用已声明的算术平均。四条印刷FESR恢复原单位后各用raw绝对.002容差、hard-midpoint截断；FF界使用ε=6×10⁻⁵，标量2mq²ε、向量ε/2。没有按输出相位更改这些量。

这些印刷输入不能全部当作精确恒等式。原文明确\(j_S=m_q(\bar uu+\bar dd)\)，并写\(m_\pi^2f_\pi^2=-\langle j_S\rangle\)；给定中心值实际有\(-Q/(m_\pi^2f_\pi^2)=.6027931\)，没有单味与味求和的二倍歧义。另以共同算术平均质量代入Eq.(2.50)，两个标量矩比打印矩低约8.48\%；RMS质量可把差别缩至约.68\%，但这只是另一味约定的代数比较，不能据接近打印数反证原作者使用了它。本轮保留打印矩与已声明FF质量规则，不偷偷``校正''凝聚或换质量规范。A2的精确输入自洽性仍有限定，其对ρ的定量影响尚未独立测定。\href{../../../evidence/FOUNDATIONS_PROOFS_20260910.md}{原输入证明与限定}

\href{../MAINLINE_SOURCE_REVIEW.md}{独立384-bit来源复核}进一步区分了两矩：n=0的raw中心差只有0.1012个当前容差，n=1则相差4.7451个容差。若额外要求n=1同时落入印刷目标和上述Eq.(2.50)中心各自±.002的盒子，两盒严格相隔超过.00549024760，无法同时满足。这个条件性不相容例子不适用于当前只用印刷目标的程序，也不是没有严格OPE余项界时对精确QCD的反证。n=1在显示阶数不含凝聚项，所以固定其它输入而只改GMOR凝聚量不能消除差别。解决时须先独立确定味电流及误差预算；本轮不修改这些物理输入。

散射双密度使用已声明的\(B(M)=100[M^2+M(M+1)/2]\)，M50时377500。它是有限正则化配方，不是物理常数。已严格证明同一有限问题的支持值随B非减、凹，并重放出\(V(B)-V(B/2)\ge.0054931289277\)。这说明正则化确实影响投影边界，不能说只是求解器预条件；但尚未证明它单独解释ρ质量或L依赖，故没有用B扫参拟合峰位。\href{../../../evidence/REGULARIZATION_VALUE_PROOF_20260911.md}{正则化证明}

原1500盘解析基线的三代表，在ρ窗口离节点存在严格η\textgreater1，并有非零T0及F1(∞)。独立完整振幅角积分重现了具体违例。这些是实际振幅问题，不能靠改纵轴或平滑消除。修补按已声明的顺序完成：

\begin{enumerate}
\def\labelenumi{\arabic{enumi}.}
\tightlist
\item
  在同一解析族追加ρ窗口及全波交错约束，再加入预定5
  MeV主波网格，总3582盘；全部原生行和电流算子保留。
\item
  对本有限族显式施加T0=0及两个FF在z=−1的双零点。双零点保证至少O(1/s)幂次；标量相应行为在原论文中是期待，未被写成已证pQCD定理。
\item
  从受控高能极限推导并加入两个线性、三个二次必要条件，随后重新取得完整严格内点、支持中心及解析primal。
\item
  在读取新相位之前固定tip、xref及二者中间截面的上边界规则，并保存完整C、ImF、R。
\end{enumerate}

这些是有明确物理理由的加强完成，不能追认为作者原始有限程序已经施加。领先高能条件严格成立只控制每个固定有限波集合的充分高能区，尚不提供全能量或无限自旋认证。新网格也只检验声明的有限样本。\href{../../../evidence/ANALYTIC_TAIL_UNITARITY_20260911.md}{尾部证明}、\href{../ENDPOINT_DESIGN_ZH.md}{端点设计}

\section{4.
完整代表的实测结果}\label{ux5b8cux6574ux4ee3ux8868ux7684ux5b9eux6d4bux7ed3ux679c}

新几何三点的中心、支持、解析核验、冻结与求值链均已完成。原有的3123盘点在独立5
MeV网格分别有98/57/59项严格违例；新增459行后，新的3582盘三点在同一185能量×3主波网格均为零严格违例。maxη=1包括阈值，不表示各处弹性。

{\def\LTcaptype{none} % do not increment counter
\begin{longtable}[]{@{}
  >{\raggedright\arraybackslash}p{(\linewidth - 8\tabcolsep) * \real{0.1579}}
  >{\raggedleft\arraybackslash}p{(\linewidth - 8\tabcolsep) * \real{0.2105}}
  >{\raggedleft\arraybackslash}p{(\linewidth - 8\tabcolsep) * \real{0.2105}}
  >{\raggedleft\arraybackslash}p{(\linewidth - 8\tabcolsep) * \real{0.2105}}
  >{\raggedleft\arraybackslash}p{(\linewidth - 8\tabcolsep) * \real{0.2105}}@{}}
\toprule\noalign{}
\begin{minipage}[b]{\linewidth}\raggedright
3582盘代表
\end{minipage} & \begin{minipage}[b]{\linewidth}\raggedleft
原生P1 90°（MeV）
\end{minipage} & \begin{minipage}[b]{\linewidth}\raggedleft
同C的5 MeV网格读数（MeV）
\end{minipage} & \begin{minipage}[b]{\linewidth}\raggedleft
原生minηP1
\end{minipage} & \begin{minipage}[b]{\linewidth}\raggedleft
原生末端δS0（°）
\end{minipage} \\
\midrule\noalign{}
\endhead
\bottomrule\noalign{}
\endlastfoot
tip & 752.363 & 747.621 & .311947 & 190.145 \\
mid & 681.596 & 680.882 & .215344 & 201.450 \\
ref & 677.890 & 679.351 & .200378 & 211.310 \\
\end{longtable}
}

\href{../fine_native_phases/phases.json}{原生数据}、\href{../fine_grid_profiles/direct_profiles.json}{网格数据}。两个90°栏目是同一个完整振幅在不同采样方式下的描述性读数，不能把它们当作两个独立物理解或严格质量误差界。

三点原生显示末端为1.1994241195
GeV；η栏是所报原生节点的最小值，不是连续能窗的严格下界。

\begin{figure}
\centering
\includegraphics[width=0.98\linewidth,height=\textheight,keepaspectratio,alt={同一IR对照与预先冻结UV三点的P1相移。论文虚线仅作输出比较，不参与物理输入或选点。}]{images/fig9.png}
\caption{同一IR对照与预先冻结UV三点的P1相移。论文虚线仅作输出比较，不参与物理输入或选点。}
\end{figure}

论文Fig.9三个可读代表的相应读数约826.62/824.30/812.53
MeV，三点跨度约14.09 MeV；本文原生跨度74.47
MeV。作者的精确代表坐标没有恢复，故这只是按公开图的描述性比较，不是宣称三点一一数学同构的排除证明。更直接的未闭合事实是：本文预选附近代表差别较大且强非弹性，S0末端也明显不同。\href{../FINE_P1_PAPER_COMPARISON.json}{图像比较记录}

原基线的五组分辨率完整计算也已交付，不能与新3582盘点拼成一个序列：

{\def\LTcaptype{none} % do not increment counter
\begin{longtable}[]{@{}lrr@{}}
\toprule\noalign{}
原基线配置(M,L) & 原生P1 90°（MeV） & minηP1 \\
\midrule\noalign{}
\endhead
\bottomrule\noalign{}
\endlastfoot
(50,8) & 835.55 & .98264 \\
(50,10) & 803.49 & .77522 \\
(50,12) & 789.08 & .62766 \\
(45,10) & 749.43 & .52188 \\
(60,10) & 773.61 & .70029 \\
\end{longtable}
}

固定M的L跨度46.47 MeV，论文约12.54 MeV，是Major差异。固定L的M跨度54.06
MeV，论文约54.35
MeV，属同量级，不能笼统称``M不收敛''。完成五组也不是M→∞定理。\href{../../sequential_reproduction_20260910/F_five_configs_verified_comparison/resolution.json}{五配置证据}

M45与M60的原生显示末端分别为1.0818369254和1.1375025560
GeV；η栏同样只总结各自实际显示点，没有把不同原生网格当作统一连续能窗。

\section{5.
Newton问题的修补与新中心复核}\label{newtonux95eeux9898ux7684ux4feeux8865ux4e0eux65b0ux4e2dux5fc3ux590dux6838}

已修复两种中心标签风险：可行恢复若改变完整C、ImF或有矩界R，不能继承中心；当前点通过中心判据后不再移动一步再保存。允许自由高能R的Schur提升，并不意味着原全变量障碍中心唯一。

八个冻结联合中心已用新规则实际重验：fine三点和原F五配置全部零Newton迭代通过，检查后未更新点；完整C、全部ImF及全部R逐值不变，原式primal与支持也同时通过。因此它们的已测ρ差异不能归因于``这几个保存点只做了一步、尚未收敛''。另有收缩历史的旧三色IR点保留独立限定，不由这八点替代。\href{../CENTER_RECHECK_RESULTS.md}{八点重验}

后续固定Watson目标曾因fresh-QR线性相对残差约1.08×10⁻⁵停止。修补允许确有下降的不精确方向继续，但中心仍要求残差≤10⁻⁶、Newton下降量及原模型曲率等≤10⁻⁸；每步另检查障碍目标实际下降。梯度/Hessian依据实际SOC密度Schur障碍，不是直接替换成四次范数的对数障碍。106项独立检查通过，核心21模块均满足文件限制，仍只有三个测试文件。\href{../../../evidence/INEXACT_NEWTON_AND_SAVED_CENTER_20260911.md}{数值证明}、\href{../verification_checked_point_descent.json}{验证记录}

固定正μ、权重及严格非空前提下，约化有限障碍有唯一内部中心，完整证明给出紧性、边界发散和自由方向Hessian正定。因此慢收敛不能当作这个有限中心根本不存在的证明；实际浮点算法仍需取得自己的最终中心及支持证据。原式支持与解析源primal分别核验，不能由一个小Newton下降量代替。

\section{6.
后续研究只用于一个明确的选择问题}\label{ux540eux7eedux7814ux7a76ux53eaux7528ux4e8eux4e00ux4e2aux660eux786eux7684ux9009ux62e9ux95eeux9898}

2403.10772增加了Watsonian选择：以旧F相位构造线性目标，在原可行域内重做完整优化。对于电流通道，\(Q=F_{old}/F_{old}^*\)，固定Q时最小化

\[1-\operatorname{Re}(Q^*S)=(1-\eta)+\eta[1-\cos2(\delta-\arg F_{old})].\]

它同时偏好弹性和与旧F的相位对齐，但必须重新取得同一个联合S/F/R点，不能把散射相位换成F相位。普通散射梯度的单调恒等式也不能直接证明这个FF驱动迭代必然全局收敛。\href{../../../evidence/WATSONIAN_DIAGNOSTIC_SCOPE_20260911.md}{预声明和证明}、\href{https://arxiv.org/html/2403.10772v4\#S2.SS5}{后续原文§2.5}

本轮只在既有3582盘模型、已冻结tip上试这一选择步骤，释放原几何坐标。后续文献改变的B-spline谱基底、额外电流、谱矩数及误差规范没有复制回2309。没有引入ρ质量、实验相移或期望峰形。第一旧目标须先完成中心、支持和解析核验，才可更新Q；初步最多比较两个完整目标更新。

\textbf{运行中状态，尚非最终物理结果：}第一固定目标已从初始J=.2595991730推进到μ=4.8828125×10⁻⁷的新标准中心J=.2808649754，上界.2929503481。支持间隙仍未达10⁻⁴，μ=2.44140625×10⁻⁷处继续定心。未完成点的相移不用于宣告ρ改善或生成第二目标。当前困难是数值路径收敛，尚无不可行证书。

\section{7.
逐项判定与保留的问题}\label{ux9010ux9879ux5224ux5b9aux4e0eux4fddux7559ux7684ux95eeux9898}

{\def\LTcaptype{none} % do not increment counter
\begin{longtable}[]{@{}
  >{\raggedright\arraybackslash}p{(\linewidth - 6\tabcolsep) * \real{0.2500}}
  >{\raggedright\arraybackslash}p{(\linewidth - 6\tabcolsep) * \real{0.2500}}
  >{\raggedright\arraybackslash}p{(\linewidth - 6\tabcolsep) * \real{0.2500}}
  >{\raggedright\arraybackslash}p{(\linewidth - 6\tabcolsep) * \real{0.2500}}@{}}
\toprule\noalign{}
\begin{minipage}[b]{\linewidth}\raggedright
论文对象
\end{minipage} & \begin{minipage}[b]{\linewidth}\raggedright
已完成
\end{minipage} & \begin{minipage}[b]{\linewidth}\raggedright
仍有差异或未完成
\end{minipage} & \begin{minipage}[b]{\linewidth}\raggedright
等级
\end{minipage} \\
\midrule\noalign{}
\endhead
\bottomrule\noalign{}
\endlastfoot
A1/A3：共同振幅与底层算子 &
完整解析有限族、交叉/投影/区间源核验、独立角积分 &
作者未公开的离散完成未唯一恢复 & 方法身份限制，不能冒称等同 \\
A2：手征/QCD输入 & 单位、味定义与印刷矩重新核对，差异显式保留 &
GMOR中心值与共同平均质量的标量矩不能同时视作精确一致 &
实质输入限定；对ρ的影响未定 \\
B2--B3、Fig.3--4 & 原1500盘基线的纯散射收缩方向见证、手征右端支持 &
完整新区域未交付，原基线手征范围偏大 & Major \\
C1--C3、Fig.5--7 & 旧三色图与新IR对照、S0零点、P1无合适ρ的控制 &
旧三色中心身份有单独限定，新全区域尚无 &
部分完成；Fig.7无ρ本身不是缺陷 \\
D1--D3：UV耦合 & 全联合严格见证、完整代表、固定IR独立UV排除 &
不能升级为所有无ρ解排除定理 & 已证结论有明确范围 \\
E1、Fig.8 & 原1500盘基线的同截面收缩支持证据 &
原基线只见轻微不对称；未复现明显不对称及完整新区域 & Major \\
E2--E3、Fig.9--10 & 同C相位/强度/谱、峰存在、新网格零严格违例 &
正确ρ的弹性、质量稳定性、S0比较 & Major \\
F1--F3、Fig.11 & 原生五配置完整计算与比较 &
L敏感性偏大；加强模型五配置未完成 & Major；M跨度单独为同量级 \\
中心保存及证明措辞 & 新代码及八点实测重验、PD/PSD等表述更正 &
不给相位误差界，不恢复成唯一QCD或连续证书 &
已修补；排除八点未通过当前有限μ中心标准这一解释 \\
\end{longtable}
}

已有充分证据的原型机缺口已经实际修补：保存点中心身份、声明网格的具体违例、已识别的T0/FF端点及领先高能系数障碍，以及原目标续算的一致性。这些修补没有完成连续全能量幺正认证。尚未完成的ρ、区域和分辨率差异仍如实保留。原论文硬约束不自动推出每个Gram谱饱和，这是可严格说明的推理范围；但原论文展示的是有限数值结果，不能据此把它改写成已被反驳的普遍定理。

更强的连续弹性阈值、全能量/无限自旋以及形状因子零点唯一性问题另列为暂缓事项。它们可能需要含额外切口的基底、零点信息或极限控制，不能未经说明改入2309输入，也不被设为完成原有限原型的新增前置条件。\href{../../../evidence/FINITE_FAMILY_INELASTIC_THRESHOLD_SCOPE_20260911.md}{暂缓阈值范围}、\href{../FOLLOWUP_SOURCE_CHECK_ZH.md}{来源核对}

最终成功标准仍是：按声明输入和事先固定的选择规则取得完整、合格的振幅，完成所请求图和比较，并准确报告实测差别。当前证据支持定性IR/UV机制及多项有限数学claims；它还不足以把Fig.3--11及正确稳健ρ的整体复现判为成功。

\end{document}
